% Table 2.1: position of this thesis relative to the closest works.
% Written by hand from paper_notes/; not generated.
% CLAUDE-DRAFT (table content)
\begin{table}[!tb]
\centering
\footnotesize
\setlength{\tabcolsep}{2.5pt}
\renewcommand{\arraystretch}{1.1}
\caption[Setting of this thesis compared with the closest works]{Setting of this thesis compared with the closest works. Pretrained model: whether the task sequence starts from a foundation model (\cref{sec:cl-medseg}). Few-shot: labelled samples per task (no: full training sets). Control task: a held-out task evaluated after the sequence to measure the loss of pretrained ability (\cref{sec:intro-approach}). \textsuperscript{a}Only for sequential fine-tuning, task-specific LoRA and shared LoRA.}
\label{tab:positioning}
\begin{tabularx}{\linewidth}{@{}L{3.9cm}L{3.3cm}L{0.6cm}L{3.4cm}L{1.6cm}>{\raggedright\arraybackslash}X@{}}
\toprule
\textbf{Work} & \textbf{Text prompt} & \textbf{3D} & \textbf{Pretrained model} & \textbf{Few-shot} & \textbf{Control task}\\
\midrule
Zhang et al., Liu et al. \mbox{}\cite{zhang2023continuallearningabdominalmultiorgan,liu2024universalextensible} & Fixed CLIP class names & yes & no (Swin UNETR) & no & no\\
MedSeqFT \mbox{}\cite{ye2025medseqft} & none & yes & yes (VoCo, UniMiSS+) & no & no\\
Chen et al. \mbox{}\cite{chen2026fewshotcontinuallearning3d} & none & yes & yes & 16, 32, 64 & no\\
MedCRP-CL \mbox{}\cite{gao2026medcrpcl} & Descriptive sentences & no & yes (CLIPSeg) & no & no\\
PA-EWC \mbox{}\cite{gao2025paewc} & Templated phrases & no & yes (CLIPSeg) & no & no\\
MedPEFT-CL \mbox{}\cite{gao2026medpeftcl} & Descriptive sentences & no & yes (CLIPSeg) & no & no\\
\midrule
This thesis & Free text, one fixed prompt per task & yes & yes (VoxTell) & 5, 10, 20\textsuperscript{a} & yes (AeroPath)\\
\bottomrule
\end{tabularx}
\end{table}
